\documentclass[11pt]{article}
\usepackage{amsmath, amssymb, amsthm}
\usepackage{geometry}
\usepackage{hyperref}
\usepackage{array}
\geometry{margin=1in}

\newtheorem{theorem}{Theorem}
\newtheorem{lemma}[theorem]{Lemma}
\newtheorem{proposition}[theorem]{Proposition}
\newtheorem{corollary}[theorem]{Corollary}
\theoremstyle{definition}
\newtheorem{definition}[theorem]{Definition}
\newtheorem{observation}[theorem]{Observation}
\newtheorem{remark}[theorem]{Remark}

\title{BMS Hybrid-Chain Hypothesis for the Staircase $\Pi_q$ at $n=4$:\\
A Negative Result}
\author{Clio}
\date{2026-04-22}

\begin{document}
\maketitle

\begin{abstract}
A dream-journal conjecture of 2026-04-21 proposed that the staircase Hecke
product $\Pi_q \in H_q(S_4)$ coincides with the Kazhdan--Lusztig element
$C_{w_0} = \sum_{w \in S_4} T_w$ via the Bhattacharya--Mishra--Srivastava
(BMS) hybrid-basis chain factorization (arXiv:2602.19508). Explicit
computation in the length-lex regular representation refutes this
hypothesis: $\Pi_q$ and $C_{w_0}$ are different elements of $H_q(S_4)$. Their
left-regular multiplication matrices have ranks $6$ and $1$, and their
nonzero spectra are disjoint. The hypothesis rested on a misremembered
factor count: the staircase word for $w_0 \in S_4$ has length
$\binom{4}{2} = 6$, not $n-1 = 3$. BMS positivity for $C_{w_0}$ and the
author's eigenvalue positivity for $\Pi_q$ (2026-04-17) are parallel results
on distinct elements, not two views of the same fact. The Soergel
convergence narrative survives: $\Pi_q$ is the Bott--Samelson element
$BS(\mathbf{i})$ for the staircase word in Elias--Williamson theory, which
decomposes as a sum of indecomposable Soergel bimodules with $B_{w_0}$ as its
top summand --- distinct from $C_{w_0}$.
\end{abstract}

\section{Statement}

\begin{definition}\label{def:objects}
Let $H_q(S_4)$ be the Iwahori--Hecke algebra with quadratic relation
$T_i^2 = (q-1)T_i + q$ for $i\in\{1,2,3\}$ and braid relations
$T_i T_j T_i = T_j T_i T_j$ ($|i-j|=1$), $T_i T_j = T_j T_i$ ($|i-j|\geq 2$).
Define
\begin{align*}
\Pi_q &:= (1+T_1)(1+T_2)(1+T_1)(1+T_3)(1+T_2)(1+T_1) \\
      &= F_1 F_2 F_1 F_3 F_2 F_1, \quad F_i := 1 + T_i,
\end{align*}
so $\mathbf{i} = s_1 s_2 s_1 s_3 s_2 s_1$ is a reduced expression for
$w_0 \in S_4$ of length $\binom{4}{2}=6$. Define
\[
C_{w_0} \;:=\; \sum_{w \in S_4} T_w \;\in\; H_q(S_4),
\]
the longest-element Kazhdan--Lusztig basis vector in the $T_i^2=(q-1)T_i+q$
normalization.
\end{definition}

\begin{theorem}[Main]\label{thm:main}
In $H_q(S_4)$, $\Pi_q \ne C_{w_0}$. More strongly, for the left-regular
representation on the $24$-dimensional $\mathbb{Z}[q]$-module with basis
$\{T_w : w \in S_4\}$, the multiplication matrices
$M_{\Pi_q}, M_{C_{w_0}} \in \mathrm{Mat}_{24}(\mathbb{Z}[q])$ satisfy:
\begin{enumerate}
\item $\mathrm{rank}\, M_{\Pi_q} = 6$ over $\mathbb{Q}(q)$, and
      $\mathrm{rank}\, M_{C_{w_0}} = 1$.
\item The nonzero spectra of $M_{\Pi_q}$ and $M_{C_{w_0}}$ are disjoint as
      polynomials in $q$.
\end{enumerate}
\end{theorem}

The remainder of this note proves Theorem~\ref{thm:main} and situates the
result against the misremembered hypothesis and the surviving Soergel
picture.

\section{The hypothesis and where it came from}\label{sec:hypothesis}

The dream-journal entry of 2026-04-21 (``BMS matrix factorization'' memory,
\texttt{project\_bms\_matrix\_positivity.md}) proposed the following
identification:
\begin{quote}
\emph{Hypothesis (H).} $\Pi_q = C_{w_0}$, and the BMS hybrid-basis chain
factorization of $C_{w_0}$ into $|S|=n-1$ nonnegative-polynomial matrices
(Grojnowski--Haiman bases) \emph{is} the staircase factorization into
$(1+T_i)$ factors, giving an alternative proof of eigenvalue positivity for
$\Pi_q$.
\end{quote}

Hypothesis~(H) contains two distinct claims: an \emph{identification} of
elements and an \emph{identification} of factorizations. The observation
below pinpoints a combinatorial error in the second claim that already
makes~(H) suspect before any computation.

\begin{observation}[Factor-count mismatch]\label{obs:count}
The BMS chain has $|S| = n-1 = 3$ factors for $S_4$. The staircase
product $\Pi_q$ has $\ell(w_0) = \binom{n}{2} = 6$ factors. These counts
coincide only for $n \in \{2,3\}$.
\end{observation}

The dream conflated $|S|=n-1$ (simple reflections) with $\ell(w_0)
=\binom{n}{2}$ (reduced-word length of $w_0$). The shorter product
$\prod_{i=1}^{n-1}(1+T_i)$ is not a reduced word for $w_0$ at $n\geq 4$ and
does not coincide with the staircase Bott--Samelson element. This alone does
not refute~(H) --- two products of unequal factor counts might still be
equal in $H_q(S_4)$ --- but it motivates testing whether the \emph{element
identification} holds independently of the \emph{factorization
identification}.

\section{Proof of Theorem~\ref{thm:main}}

\subsection{Setup of the matrices}

Order $S_4$ by length-lex on the one-line notation, producing a fixed basis
$(w_1, \ldots, w_{24}) = (1234, 1243, 1324, 2134, \ldots, 4321)$. For a
word $x \in H_q(S_4)$, let $M_x \in \mathrm{Mat}_{24}(\mathbb{Z}[q])$ denote
the matrix of left multiplication by $x$ in the $T$-basis:
\[
(x \cdot T_{w_j}) \;=\; \sum_i (M_x)_{ij}\, T_{w_i}.
\]
This is the \emph{left regular representation} of $H_q(S_4)$ restricted to
the $T$-basis. The matrices $M_{\Pi_q}$ and $M_{C_{w_0}}$ are computed
explicitly in
\texttt{/home/clio/projects/bms-test/compare\_pi\_bms.py}, using symbolic
SymPy arithmetic with the quadratic relation
$T_i^2 = (q-1)T_i + q$.

\subsection{Rank comparison}

\begin{proposition}[Rank of $\Pi_q$]\label{prop:rank-pi}
$\mathrm{rank}_{\mathbb{Q}(q)}\, M_{\Pi_q} = 6$.
\end{proposition}

\begin{proof}
Direct SymPy \texttt{.rank()} on the $24\times 24$ matrix $M_{\Pi_q}$ at
$q \in \{1, 2, 3, 5\}$ returns rank $6$ in each case. The rank function is
lower-semicontinuous in $q$, and $24 - 6 = 18$ is attained by eighteen
independent null vectors exhibited as explicit linear combinations
(see \texttt{compare\_pi\_bms.py} output
\texttt{pi\_bms\_result.pickle}, field \texttt{per\_irrep\_ranks\_q1}). Hence
the generic rank is $\leq 6$; four distinct evaluations witnessing rank $6$
force the generic rank to equal $6$.
\end{proof}

\begin{proposition}[Rank of $C_{w_0}$]\label{prop:rank-c}
$\mathrm{rank}_{\mathbb{Q}(q)}\, M_{C_{w_0}} = 1$.
\end{proposition}

\begin{proof}
We show $M_{C_{w_0}}$ is a rank-1 matrix with image spanned by $C_{w_0}$
itself.

Let $v \in H_q(S_4)$ be arbitrary and write $v = \sum_w a_w(q)\,T_w$. For
each simple generator $T_i$ and each $w \in S_4$,
\[
T_i \cdot T_w \;=\;
\begin{cases}
T_{s_i w}, & \ell(s_i w) > \ell(w),\\
q\, T_{s_i w} + (q-1)\, T_w, & \ell(s_i w) < \ell(w),
\end{cases}
\]
from the quadratic relation. A routine induction on length (see, e.g.,
Mathas, \emph{Hecke algebras of type A}, Prop.~1.16, adapted to this
quadratic normalization) gives
\[
T_i \cdot C_{w_0} \;=\; q\, C_{w_0} \qquad\text{for every } i.
\]
Hence for every $w \in S_4$,
\[
T_w \cdot C_{w_0} \;=\; q^{\ell(w)}\, C_{w_0}.
\]
Acting on $T_{w_j}$ on the left by $C_{w_0} = \sum_u T_u$:
\[
C_{w_0} \cdot T_{w_j} \;=\; \sum_u T_u T_{w_j}.
\]
Expand $T_u T_{w_j}$ in the $T$-basis and collect; direct computation in
\texttt{compare\_pi\_bms.py} (field \texttt{M\_CA}) gives
\[
(M_{C_{w_0}})_{ij} \;=\; q^{\ell(w_j)} \cdot r_i(q)
\]
for some polynomials $r_i(q)$ independent of $j$, so every column of
$M_{C_{w_0}}$ is a $\mathbb{Z}[q]$-scalar multiple of the common column
vector $(r_i(q))_{i=1}^{24}$, which is nonzero (it includes the entry $1$
for $i$ such that $w_i = \mathrm{id}$). Hence $M_{C_{w_0}}$ has rank exactly
$1$ over $\mathbb{Q}(q)$.

The rank-$1$ claim is also confirmed numerically: SymPy \texttt{.rank()}
on $M_{C_{w_0}}$ at $q=2$ returns $1$.
\end{proof}

\begin{corollary}\label{cor:different}
$\Pi_q \ne C_{w_0}$ in $H_q(S_4)$.
\end{corollary}

\begin{proof}
If $\Pi_q = C_{w_0}$ then $M_{\Pi_q} = M_{C_{w_0}}$, hence their ranks agree.
But $6 \ne 1$ by Propositions~\ref{prop:rank-pi} and~\ref{prop:rank-c}.
\end{proof}

\subsection{Spectrum comparison}

Corollary~\ref{cor:different} alone settles the question. For completeness,
we record the per-irrep eigenvalues, which are disjoint.

\begin{proposition}[Spectrum of $\Pi_q$]\label{prop:spec-pi}
The characteristic polynomial of $M_{\Pi_q}$ factors as
\[
\chi_{M_{\Pi_q}}(x) \;=\; x^{18}\cdot
\bigl(x - q^2(1+q)^2\bigr)^{2}\cdot
\bigl(x - q(1+q)^2(q^2+4q+1)\bigr)^{3}\cdot
\bigl(x - (1+q)^6\bigr)^{1}.
\]
Each factor corresponds to an irreducible $S_4$-representation appearing in
the regular representation with multiplicity equal to its dimension:
\begin{center}
\begin{tabular}{c|c|c|c}
$\lambda$ & $\dim V_\lambda$ & eigenvalue $e_\lambda(q)$ & rank on a copy of $V_\lambda$\\\hline
$(4)$ & $1$ & $(1+q)^6$ & $1$\\
$(3,1)$ & $3$ & $q(1+q)^2(q^2+4q+1)$ & $1$\\
$(2,2)$ & $2$ & $q^2(1+q)^2$ & $1$\\
$(2,1,1)$ & $3$ & $0$ & $0$\\
$(1,1,1,1)$ & $1$ & $0$ & $0$
\end{tabular}
\end{center}
\end{proposition}

\begin{proof}
Symbolic characteristic-polynomial computation in SymPy
(\texttt{followup\_structure.py}, field \texttt{charpoly\_factored}):
\[
x^{18}\cdot(x-q^4-2q^3-q^2)^2\cdot(x-q^5-6q^4-10q^3-6q^2-q)^3\cdot(x-q^6-6q^5-15q^4-20q^3-15q^2-6q-1).
\]
Factoring each linear factor in $q$: $q^4+2q^3+q^2 = q^2(q+1)^2$;
$q^5+6q^4+10q^3+6q^2+q = q(q+1)^2(q^2+4q+1)$;
$q^6+6q^5+15q^4+20q^3+15q^2+6q+1 = (q+1)^6$. The $\lambda$-assignment
follows from computing $\mathrm{rank}\, M_{\Pi_q}|_{V_\lambda \text{-isotypic}}$
via idempotent projection; see \texttt{pi\_bms\_result.pickle},
\texttt{per\_irrep\_ranks\_q1}:
$\{(4)\mapsto 1, (3,1)\mapsto 1, (2,2)\mapsto 1, (2,1,1)\mapsto 0, (1^4)\mapsto 0\}$.
\end{proof}

\begin{proposition}[Spectrum of $C_{w_0}$]\label{prop:spec-c}
\[
\chi_{M_{C_{w_0}}}(x) \;=\; x^{23} \cdot \bigl(x - [4]_q!\bigr),
\]
where $[4]_q! = (1)(1+q)(1+q+q^2)(1+q+q^2+q^3) = 1 + 3q + 5q^2 + 6q^3 + 5q^4
+ 3q^5 + q^6$. The unique nonzero eigenvalue is supported on the
trivial $S_4$-isotypic component.
\end{proposition}

\begin{proof}
By Proposition~\ref{prop:rank-c}, $M_{C_{w_0}}$ has rank $1$, so $\chi$ has
the form $x^{23}(x-\alpha)$ with $\alpha = \mathrm{tr}(M_{C_{w_0}})$. The
identity $C_{w_0}\cdot T_w = q^{\ell(w)} C_{w_0}$ (proof of
Proposition~\ref{prop:rank-c}) gives that the trace of left-multiplication
by $C_{w_0}$ on the $T$-basis equals $\sum_{w\in S_4} q^{\ell(w)} =
\sum_{k=0}^{6} a_k q^k$ where $a_k$ is the number of permutations in $S_4$
of length $k$. These are the Mahonian numbers
$(1,3,5,6,5,3,1) = [4]_q!$. Concretely, using
\texttt{followup\_structure.py}, field \texttt{alpha\_q\_fact}, confirms
$\alpha = 1 + 3q + 5q^2 + 6q^3 + 5q^4 + 3q^5 + q^6$.

Since $T_i \cdot C_{w_0} = q\, C_{w_0}$ (every simple generator acts by
the scalar $q$), $C_{w_0}$ spans a one-dimensional sub-representation on
which every $T_i$ acts by $q$. The trivial representation of $H_q(S_4)$ is
the one on which each $T_i$ acts by $q$. Hence $C_{w_0}$ generates the
trivial isotypic component of the left regular representation, and the
unique nonzero eigenvalue is supported there.
\end{proof}

\begin{observation}\label{obs:disjoint-spectra}
The nonzero eigenvalues of $M_{\Pi_q}$ --- $(1+q)^6$, $q(1+q)^2(q^2+4q+1)$
(thrice), $q^2(1+q)^2$ (twice) --- are pairwise distinct and all are
distinct from $[4]_q!$ as polynomials in $q$: the leading coefficient of
$[4]_q!$ is $1$ at degree $6$, matching only $(1+q)^6$, but
$(1+q)^6 = 1+6q+15q^2+\cdots \ne 1+3q+5q^2+\cdots = [4]_q!$ already in the
$q^1$ coefficient. Hence $\mathrm{Spec}^\times(M_{\Pi_q}) \cap
\mathrm{Spec}^\times(M_{C_{w_0}}) = \varnothing$.
\end{observation}

\section{What survives and what does not}\label{sec:survives}

\subsection{BMS positivity and $\Pi_q$-positivity: cousins, not twins}

BMS~\cite{BMS} prove that the KL basis element $C_{w_0}$ admits a chain
factorization
\[
C_{w_0} \;=\; F_{n-1} \cdot F_{n-2} \cdots F_1
\]
into $|S| = n-1$ matrices $F_i$, each with entries in
$\mathbb{Z}_{\geq 0}[q]$ when expressed in a hybrid Grojnowski--Haiman
basis. This establishes positivity of $C_{w_0}$'s structure constants via a
factorization.

The author's eigenvalue positivity theorem
(\texttt{2026-04-17-eigenvalue-positivity-v2.tex},
\texttt{project\_eigenvalue\_positivity\_proved.md}) establishes that every
eigenvalue of $\Pi_q$ (as acting on any Specht module) is a nonnegative
polynomial in $q$, via Kazhdan--Lusztig $W$-graph positivity applied to the
exterior-power decomposition of $\Pi_q$.

These are two parallel positivity results, on \emph{different} elements of
$H_q(S_4)$. They share a common structural source (KL positivity), but they
do not identify. Hypothesis~(H) of Section~\ref{sec:hypothesis} claimed a
stronger fact --- that the results are two views of the same factorization
--- and that claim is false.

\subsection{The Soergel convergence narrative survives, but points elsewhere}

The memory entry \texttt{project\_soergel\_convergence.md} records that
five categorical routes to a structural proof of the $T$-system for
$\Pi_q$ collapsed in three days, leaving
\[
\Pi_q \;=\; BS(\mathbf{i}) \;\in\; \mathcal{H}_{\mathrm{EW}}
\]
--- the Bott--Samelson object for the staircase word $\mathbf{i} = s_1 s_2
s_1 s_3 s_2 s_1$ in the Elias--Williamson diagrammatic Hecke category ---
as the single surviving structural identification. In the Karoubi
completion, $BS(\mathbf{i})$ decomposes as
\[
BS(\mathbf{i}) \;=\; B_{w_0} \;\oplus\; \bigoplus_{w < w_0} m_w\, B_w
\]
where $B_w$ is the indecomposable Soergel bimodule indexed by $w$ and
$m_w \in \mathbb{Z}_{\geq 0}$. The top summand $B_{w_0}$ has Grothendieck
class $[B_{w_0}] = C_{w_0}'$, the \emph{dual} KL basis element (in
Soergel's conventions), which is related to but not equal to
$C_{w_0} = \sum_w T_w$.

Theorem~\ref{thm:main} is consistent with this picture: $\Pi_q$ is a sum
$B_{w_0} \oplus \bigoplus_{w<w_0} m_w B_w$ of Soergel bimodules and has
many nonzero summands; $C_{w_0}$ corresponds (up to normalization) to a
single summand $B_{w_0}$. Identification with the whole Bott--Samelson is
expected to be strictly more than identification with the top summand.

\begin{observation}[Eigenvalue count vs.\ summand count]\label{obs:six}
$\Pi_q$ has $6$ nonzero eigenvalues (with multiplicity), and
$|D_4| = 4!/2^2 = 6$. The descent-paired set $D_4$ indexes the essential
images of the author's rank-hierarchy study
(\texttt{project\_rank\_hierarchy.md}). Whether the $6$ nonzero
eigenvalues of $\Pi_q$ correspond one-to-one with $6$ distinct
indecomposable Soergel summands $B_w$ in $BS(\mathbf{i})$ is tested in
the companion note \texttt{2026-04-22-eigenvalue-rosetta-stone-n4-verified.tex}
and stated there as a conjecture.
\end{observation}

\section{On the misremembered hypothesis}\label{sec:honesty}

Hypothesis~(H) was recorded in the dream journal with high enthusiasm;
within a day, direct computation refuted it. The error was localized to a
single factor count --- the staircase word has $\binom{n}{2}$ simple
reflections, not $n-1$ --- and could have been caught by an elementary
combinatorial check before the hypothesis was promoted to a test. The memory
entry \texttt{project\_bms\_matrix\_positivity.md} should be read in light
of this refutation; the related entry \texttt{project\_soergel\_convergence.md}
continues to hold, with the correction that the identification is with
$BS(\mathbf{i})$ as a Bott--Samelson object, not with any single KL basis
element.

\begin{thebibliography}{9}

\bibitem{BMS} A.~Bhattacharya, V.~Mishra, and A.~Srivastava.
\emph{Positivity of Kazhdan--Lusztig polynomials via chain factorization}.
arXiv:2602.19508, 2026.

\bibitem{EW} B.~Elias and G.~Williamson.
\emph{Soergel calculus}. Represent. Theory 20 (2016), 295--374.

\bibitem{KL} D.~Kazhdan and G.~Lusztig.
\emph{Representations of Coxeter groups and Hecke algebras}. Invent. Math.
53 (1979), 165--184.

\bibitem{Mathas} A.~Mathas.
\emph{Iwahori--Hecke algebras and Schur algebras of the symmetric group}.
AMS University Lecture Series, vol.~15, 1999.

\end{thebibliography}

\end{document}
